\documentclass{ctexart}
\usepackage{geometry}
\usepackage{listings}
\geometry{a4paper, margin=1in}
\lstset{basicstyle=\ttfamily\small, breaklines=true, frame=single}
\title{2\_7 修改记录}
\date{2026-02-07}
\begin{document}
\maketitle

\section{train\_balanced\_sampling\_sep\_mlp.py 改为 HGAT + 双MLP(源/目标域)}
\subsection{修改内容概述}
\begin{itemize}
  \item 删除旧的 Path/CNN 训练逻辑与相关辅助函数，仅保留 dataset.pkl + HGAT 的训练流程。
  \item 新增数值特征规范化与 CellDelayDataset，直接读取 dataset.pkl 中的 DataFrame。
  \item 使用 HGAT 编码 SPICE 图，加入 src/tgt 双 MLP head（分别用于源域/目标域）。
  \item 训练入口改为 train\_balanced\_sep\_mlp，验证仅在目标域 val 上进行。
\end{itemize}

\subsection{修改前后对比（数据读取）}
\textbf{修改前：} 设计图 pkl + path/CNN 数据
\begin{lstlisting}
train_dataset_7 = load_data(...)
train_dataset_130 = load_data(...)
val_dataset = load_data(...)
\end{lstlisting}

\textbf{修改后：} dataset.pkl 读取 df
\begin{lstlisting}
obj = load_dataset_pkl(data_dir, options.dataset_pkl_name)
df_tgt_train = obj.get("tgt_train_df")
df_src = obj.get("src_df")
\end{lstlisting}

\subsection{修改前后对比（模型/训练）}
\textbf{修改前：} PathConv + CNN + PathModel
\begin{lstlisting}
cur_label_hats = model(graph, nodes, eids, targets, level_id, path_map)
\end{lstlisting}

\textbf{修改后：} HGAT + 双MLP回归（src/tgt）
\begin{lstlisting}
enc = HGATDesignEncoder(...)
model = SepCellDelayRegressor(...)
pred_tgt = model(xb_tgt, zb_tgt, node="tgt")
pred_src = model(xb_src, zb_src, node="src")
\end{lstlisting}

\subsection{参数与训练入口调整}
\textbf{修改前：}
\begin{lstlisting}
train(options, seed)
\end{lstlisting}

\textbf{修改后：}
\begin{lstlisting}
train_balanced_sep_mlp(options, seed)
\end{lstlisting}

\subsection{options.py 说明更新}
\textbf{修改前：}
\begin{lstlisting}
--in_dim ... "the dimension of the input feature"
--data_save_path ... "directory that contains the dataset"
\end{lstlisting}

\textbf{修改后：}
\begin{lstlisting}
--in_dim ... "must equal len(NUMERIC_COLS)=20"
--data_save_path ... "directory containing dataset.pkl"
\end{lstlisting}

\end{document}
